\section{U}

% üben (to practice)
\begin{verbentry}{
  style = {default},
  toc = {üben (to practice)},
  name = {üben},
  english = {to practice},
  type = {regular},
  auxiliary = {haben}
}
 
    
     \vspace{5pt}

    \hrule
    
    \vspace{5pt}

  \begin{tabular}{rl}
     \multicolumn{2}{l}{\textbf{PRÄSENS} }\\
    ich & \textbf{übe}\\
    du & \textbf{übst}\\
    er/sie/es & \textbf{übt}\\
    wir & \textbf{üben}\\
    ihr & \textbf{übt}\\
    sie/Sie & \textbf{üben}\\
  \end{tabular}
  \hfill
     \begin{tabular}{rl}
    \multicolumn{2}{l}{\textbf{PRÄTERITUM} }\\
    ich & \textbf{übte}\\
    du & \textbf{übtest}\\
    er/sie/es & \textbf{übte}\\
    wir & \textbf{übten}\\
    ihr & \textbf{übtet}\\
    sie/Sie & \textbf{übten}\\
  \end{tabular}
  \hfill
  \begin{tabular}{rl}
     \multicolumn{2}{l}{\textbf{IMPERATIV} }\\
   & \textbf{üb(e) (du)!}\\
   & \textbf{üben wir!}\\
   & \textbf{übt (ihr)!}\\
   & \textbf{üben Sie!}\\
   & \vspace{10pt}
  \end{tabular}

    \vspace{10pt}

 \begin{tabular}{rl}
     \multicolumn{2}{l}{\textbf{PERFEKT}}\\
   & haben + \textbf{geübt}\\
  \end{tabular}
\hfill
   \begin{tabular}{rl}
    \multicolumn{2}{l}{\textbf{FUTURE I} }\\
   & werde + \textbf{üben}\\
  \end{tabular}
\hfill
   \begin{tabular}{rl}
      \multicolumn{2}{l}{\textbf{PLUSQUAMPERFEKT}}\\
   & hatten + \textbf{geübt}\\
  \end{tabular}

   \vspace{10pt}

   \begin{tabular}{lll}
      \multicolumn{3}{l}{\textbf{MIT PRÄPOSITIONEN}}\\
& \textbf{üben für} + Akk & (to practice for)\\
& \textbf{üben mit} + Dat & (to practice with)\\
\end{tabular}
  \hfill
   \begin{tabular}{lll}
      \multicolumn{3}{l}{\textbf{TRENNBARE VERBEN}}\\
 & \textbf{ein$|$üben} & (to rehearse / practice)\\
 & \textbf{aus$|$üben} & (to exercise / carry out)\\
\vspace{10pt}
  \end{tabular}

  \vspace{-5pt}
\end{verbentry}

% überlegen (to think about)
\begin{verbentry}{
  style = {acc},
  toc = {überlegen (to think about)},
  name = {überlegen},
  english = {to think about},
  type = {irregular, + Akkusativ},
  auxiliary = {haben}
}
 
    
     \vspace{5pt}

    \hrule
    
    \vspace{5pt}

  \begin{tabular}{rl}
     \multicolumn{2}{l}{\textbf{PRÄSENS} }\\
    ich & \textbf{überlege}\\
    du & \textbf{überlegst}\\
    er/sie/es & \textbf{überlegt}\\
    wir & \textbf{überlegen}\\
    ihr & \textbf{überlegt}\\
    sie/Sie & \textbf{überlegen}\\
  \end{tabular}
  \hfill
     \begin{tabular}{rl}
    \multicolumn{2}{l}{\textbf{PRÄTERITUM} }\\
    ich & \textbf{überlegte}\\
    du & \textbf{überlegtest}\\
    er/sie/es & \textbf{überlegte}\\
    wir & \textbf{überlegten}\\
    ihr & \textbf{überlegtet}\\
    sie/Sie & \textbf{überlegten}\\
  \end{tabular}
  \hfill
  \begin{tabular}{rl}
     \multicolumn{2}{l}{\textbf{IMPERATIV} }\\
   & \textbf{überleg(e) (du)!}\\
   & \textbf{überlegen wir!}\\
   & \textbf{überlegt (ihr)!}\\
   & \textbf{überlegen Sie!}\\
   & \vspace{10pt}
  \end{tabular}

    \vspace{10pt}

 \begin{tabular}{rl}
     \multicolumn{2}{l}{\textbf{PERFEKT}}\\
   & haben + \textbf{überlegt}\\
  \end{tabular}
\hfill
   \begin{tabular}{rl}
   
    \multicolumn{2}{l}{\textbf{FUTURE I} }\\
  &  werden + \textbf{überlegen}\\
 
  \end{tabular}
\hfill
   \begin{tabular}{rl}
      \multicolumn{2}{l}{\textbf{PLUSQUAMPERFEKT}}\\
  &  hatten + \textbf{überlegt}\\
  \end{tabular}
  
  
\vspace{10pt}

\begin{tabular}{lll}
\multicolumn{3}{l}{\textbf{MIT PRÄPOSITIONEN}}\\
& \textbf{---} & \\
\end{tabular}
\hfill
\begin{tabular}{lll}
\multicolumn{3}{l}{\textbf{TRENNBARE VERBEN}}\\
& \textbf{---} & \\
\end{tabular}
\vspace{-5pt}
\end{verbentry}

% übernehmen (to take over / to assume)
\begin{verbentry}{
  style = {default},
  toc = {übernehmen (to take over / to assume)},
  name = {übernehmen},
  english = {to take over / to assume},
  type = {irregular},
  auxiliary = {haben}
}
 
    
     \vspace{5pt}

    \hrule
    
    \vspace{5pt}

  \begin{tabular}{rl}
     \multicolumn{2}{l}{\textbf{PRÄSENS} }\\
    ich & \textbf{übernehme}\\
    du & \textbf{übernimmst}\\
    er/sie/es & \textbf{übernimmt}\\
    wir & \textbf{übernehmen}\\
    ihr & \textbf{übernehmt}\\
    sie/Sie & \textbf{übernehmen}\\
  \end{tabular}
  \hfill
     \begin{tabular}{rl}
    \multicolumn{2}{l}{\textbf{PRÄTERITUM} }\\
    ich & \textbf{übernahm}\\
    du & \textbf{übernahmst}\\
    er/sie/es & \textbf{übernahm}\\
    wir & \textbf{übernahmen}\\
    ihr & \textbf{übernahmt}\\
    sie/Sie & \textbf{übernahmen}\\
  \end{tabular}
  \hfill
  \begin{tabular}{rl}
     \multicolumn{2}{l}{\textbf{IMPERATIV} }\\
   & \textbf{übernimm (du)!}\\
   & \textbf{übernehmen wir!}\\
   & \textbf{übernehmt (ihr)!}\\
   & \textbf{übernehmen Sie!}\\
   & \vspace{10pt}
  \end{tabular}

    \vspace{10pt}

 \begin{tabular}{rl}
     \multicolumn{2}{l}{\textbf{PERFEKT}}\\
   & haben + \textbf{übernommen}\\
  \end{tabular}
\hfill
   \begin{tabular}{rl}
    \multicolumn{2}{l}{\textbf{FUTURE I} }\\
   & werden + \textbf{übernehmen}\\
  \end{tabular}
\hfill
   \begin{tabular}{rl}
      \multicolumn{2}{l}{\textbf{PLUSQUAMPERFEKT}}\\
   & hatten + \textbf{übernommen}\\
  \end{tabular}

   \vspace{10pt}

   \begin{tabular}{lll}
      \multicolumn{3}{l}{\textbf{MIT PRÄPOSITIONEN}}\\
& \textbf{übernehmen von} + Dat & (to take over from)\\
& \textbf{übernehmen für} + Akk & (to assume responsibility for)\\
\end{tabular}
  \hfill
   \begin{tabular}{lll}
\multicolumn{3}{l}{\textbf{TRENNBARE VERBEN}}\\
& \textbf{---} & \\
\end{tabular}

  \vspace{-5pt}
\end{verbentry}

% überweisen (to transfer money / refer)
\begin{verbentry}{
  style = {default},
  toc = {überweisen (to transfer money / refer)},
  name = {überweisen},
  english = {to transfer money / refer},
  type = {regular},
  auxiliary = {haben}
}


\vspace{5pt}
\hrule
\vspace{5pt}

\begin{tabular}{rl}
\multicolumn{2}{l}{\textbf{PRÄSENS}}\\
ich & \textbf{überweise}\\
du & \textbf{überweist}\\
er/sie/es & \textbf{überweist}\\
wir & \textbf{überweisen}\\
ihr & \textbf{überweist}\\
sie/Sie & \textbf{überweisen}\\
\end{tabular}
\hfill
\begin{tabular}{rl}
\multicolumn{2}{l}{\textbf{PRÄTERITUM}}\\
ich & \textbf{überwies}\\
du & \textbf{überwiesest}\\
er/sie/es & \textbf{überwies}\\
wir & \textbf{überwiesen}\\
ihr & \textbf{überwieset}\\
sie/Sie & \textbf{überwiesen}\\
\end{tabular}
\hfill
\begin{tabular}{rl}
\multicolumn{2}{l}{\textbf{IMPERATIV}}\\
& \textbf{überweise (du)!}\\
& \textbf{überweisen wir!}\\
& \textbf{überweist (ihr)!}\\
& \textbf{überweisen Sie!}\\
\end{tabular}

\vspace{10pt}

\begin{tabular}{rl}
\multicolumn{2}{l}{\textbf{PERFEKT}}\\
& haben + \textbf{überwiesen}\\
\end{tabular}
\hfill
\begin{tabular}{rl}
\multicolumn{2}{l}{\textbf{FUTURE I}}\\
& werden + \textbf{überweisen}\\
\end{tabular}
\hfill
\begin{tabular}{rl}
\multicolumn{2}{l}{\textbf{PLUSQUAMPERFEKT}}\\
& hatten + \textbf{überwiesen}\\
\end{tabular}

\vspace{10pt}

\begin{tabular}{lll}
\multicolumn{3}{l}{\textbf{MIT PRÄPOSITIONEN}}\\
& \textbf{überweisen an} + Akk & (to transfer to / refer to)\\
\end{tabular}
\hfill
\begin{tabular}{lll}
\multicolumn{3}{l}{\textbf{TRENNBARE VERBEN}}\\
\vspace{10pt}\\
\end{tabular}
\vspace{-5pt}
\end{verbentry}

% umarmen (to hug)
\begin{verbentry}{
  style = {acc},
  toc = {umarmen (to hug)},
  name = {umarmen},
  english = {to hug},
  type = {regular, + Akkusativ},
  auxiliary = {haben}
}


\vspace{5pt}
\hrule
\vspace{5pt}

\begin{tabular}{rl}
\multicolumn{2}{l}{\textbf{PRÄSENS}}\\
ich & \textbf{umarme}\\
du & \textbf{umarmst}\\
er/sie/es & \textbf{umarmt}\\
wir & \textbf{umarmen}\\
ihr & \textbf{umarmt}\\
sie/Sie & \textbf{umarmen}\\
\end{tabular}
\hfill
\begin{tabular}{rl}
\multicolumn{2}{l}{\textbf{PRÄTERITUM}}\\
ich & \textbf{umarmte}\\
du & \textbf{umarmtest}\\
er/sie/es & \textbf{umarmte}\\
wir & \textbf{umarmten}\\
ihr & \textbf{umarmtet}\\
sie/Sie & \textbf{umarmten}\\
\end{tabular}
\hfill
\begin{tabular}{rl}
\multicolumn{2}{l}{\textbf{IMPERATIV}}\\
& \textbf{umarme (du)!}\\
& \textbf{umarmen wir!}\\
& \textbf{umarmt (ihr)!}\\
& \textbf{umarmen Sie!}\\
\end{tabular}

\vspace{10pt}

\begin{tabular}{rl}
\multicolumn{2}{l}{\textbf{PERFEKT}}\\
& haben + \textbf{umarmt}\\
\end{tabular}
\hfill
\begin{tabular}{rl}
\multicolumn{2}{l}{\textbf{FUTURE I}}\\
& werden + \textbf{umarmen}\\
\end{tabular}
\hfill
\begin{tabular}{rl}
\multicolumn{2}{l}{\textbf{PLUSQUAMPERFEKT}}\\
& hatten + \textbf{umarmt}\\
\end{tabular}

\vspace{10pt}

\begin{tabular}{lll}
\multicolumn{3}{l}{\textbf{MIT PRÄPOSITIONEN}}\\
& \textbf{---} & \\
\end{tabular}
\hfill
\begin{tabular}{lll}
\multicolumn{3}{l}{\textbf{TRENNBARE VERBEN}}\\
& \textbf{---} & \\
\end{tabular}

\vspace{-5pt}
\end{verbentry}

% unterhalten (to entertain)
\begin{verbentry}{
  style = {acc},
  toc = {unterhalten (to entertain)},
  name = {unterhalten},
  english = {to entertain},
  type = {irregular, + Akkusativ},
  auxiliary = {haben}
}


\vspace{5pt}
\hrule
\vspace{5pt}

\begin{tabular}{rl}
\multicolumn{2}{l}{\textbf{PRÄSENS}}\\
ich & \textbf{unterhalte}\\
du & \textbf{unterhältst}\\
er/sie/es & \textbf{unterhält}\\
wir & \textbf{unterhalten}\\
ihr & \textbf{unterhaltet}\\
sie/Sie & \textbf{unterhalten}\\
\end{tabular}
\hfill
\begin{tabular}{rl}
\multicolumn{2}{l}{\textbf{PRÄTERITUM}}\\
ich & \textbf{unterhielt}\\
du & \textbf{unterhieltest}\\
er/sie/es & \textbf{unterhielt}\\
wir & \textbf{unterhielten}\\
ihr & \textbf{unterhieltet}\\
sie/Sie & \textbf{unterhielten}\\
\end{tabular}
\hfill
\begin{tabular}{rl}
\multicolumn{2}{l}{\textbf{IMPERATIV}}\\
& \textbf{unterhalte (du)!}\\
& \textbf{unterhalten wir!}\\
& \textbf{unterhaltet (ihr)!}\\
& \textbf{unterhalten Sie!}\\
\end{tabular}

\vspace{10pt}

\begin{tabular}{rl}
\multicolumn{2}{l}{\textbf{PERFEKT}}\\
& haben + \textbf{unterhalten}\\
\end{tabular}
\hfill
\begin{tabular}{rl}
\multicolumn{2}{l}{\textbf{FUTURE I}}\\
& werden + \textbf{unterhalten}\\
\end{tabular}
\hfill
\begin{tabular}{rl}
\multicolumn{2}{l}{\textbf{PLUSQUAMPERFEKT}}\\
& hatten + \textbf{unterhalten}\\
\end{tabular}

\vspace{10pt}

\begin{tabular}{lll}
\multicolumn{3}{l}{\textbf{MIT PRÄPOSITIONEN}}\\
& \textbf{---} & \\
\end{tabular}
\hfill
\begin{tabular}{lll}
\multicolumn{3}{l}{\textbf{TRENNBARE VERBEN}}\\
& \textbf{---} & \\
\end{tabular}

\vspace{-5pt}
\end{verbentry}

% unternehmen (to undertake / to do)
\begin{verbentry}{
  style = {default},
  toc = {unternehmen (to undertake / to do)},
  name = {unternehmen},
  english = {to undertake / to do},
  type = {irregular},
  auxiliary = {haben}
}
 
    
     \vspace{5pt}

    \hrule
    
    \vspace{5pt}

  \begin{tabular}{rl}
     \multicolumn{2}{l}{\textbf{PRÄSENS} }\\
    ich & \textbf{unternehme}\\
    du & \textbf{unternimmst}\\
    er/sie/es & \textbf{unternimmt}\\
    wir & \textbf{unternehmen}\\
    ihr & \textbf{unternehmt}\\
    sie/Sie & \textbf{unternehmen}\\
  \end{tabular}
  \hfill
     \begin{tabular}{rl}
    \multicolumn{2}{l}{\textbf{PRÄTERITUM} }\\
    ich & \textbf{unternahm}\\
    du & \textbf{unternahmst}\\
    er/sie/es & \textbf{unternahm}\\
    wir & \textbf{unternahmen}\\
    ihr & \textbf{unternahmt}\\
    sie/Sie & \textbf{unternahmen}\\
  \end{tabular}
  \hfill
  \begin{tabular}{rl}
     \multicolumn{2}{l}{\textbf{IMPERATIV} }\\
   & \textbf{unternimm (du)!}\\
   & \textbf{unternehmen wir!}\\
   & \textbf{unternehmt (ihr)!}\\
   & \textbf{unternehmen Sie!}\\
   & \vspace{10pt}
  \end{tabular}

    \vspace{10pt}

 \begin{tabular}{rl}
     \multicolumn{2}{l}{\textbf{PERFEKT}}\\
   & haben + \textbf{unternommen}\\
  \end{tabular}
\hfill
   \begin{tabular}{rl}
    \multicolumn{2}{l}{\textbf{FUTURE I} }\\
   & werden + \textbf{unternehmen}\\
  \end{tabular}
\hfill
   \begin{tabular}{rl}
      \multicolumn{2}{l}{\textbf{PLUSQUAMPERFEKT}}\\
   & hatten + \textbf{unternommen}\\
  \end{tabular}

   \vspace{10pt}

   \begin{tabular}{lll}
      \multicolumn{3}{l}{\textbf{MIT PRÄPOSITIONEN}}\\
& \textbf{unternehmen gegen} + Akk & (to undertake against)\\
& \textbf{unternehmen mit} + Dat & (to do with)\\
\end{tabular}
  \hfill
   \begin{tabular}{lll}
\multicolumn{3}{l}{\textbf{TRENNBARE VERBEN}}\\
& \textbf{---} & \\
\end{tabular}

  \vspace{-5pt}
\end{verbentry}

% unterscheiden (to distinguish)
\begin{verbentry}{
  style = {default},
  toc = {unterscheiden (to distinguish)},
  name = {unterscheiden},
  english = {to distinguish},
  type = {irregular},
  auxiliary = {haben}
}
 

\vspace{5pt}
\hrule
\vspace{5pt}

\begin{tabular}{rl}
\multicolumn{2}{l}{\textbf{PRÄSENS}}\\
ich & \textbf{unterscheide}\\
du & \textbf{unterscheidest}\\
er/sie/es & \textbf{unterscheidet}\\
wir & \textbf{unterscheiden}\\
ihr & \textbf{unterscheidet}\\
sie/Sie & \textbf{unterscheiden}\\
\end{tabular}
\hfill
\begin{tabular}{rl}
\multicolumn{2}{l}{\textbf{PRÄTERITUM}}\\
ich & \textbf{unterschied}\\
du & \textbf{unterschiedst}\\
er/sie/es & \textbf{unterschied}\\
wir & \textbf{unterschieden}\\
ihr & \textbf{unterschiedet}\\
sie/Sie & \textbf{unterschieden}\\
\end{tabular}
\hfill
\begin{tabular}{rl}
\multicolumn{2}{l}{\textbf{IMPERATIV}}\\
& \textbf{unterscheide (du)!}\\
& \textbf{unterscheiden wir!}\\
& \textbf{unterscheidet (ihr)!}\\
& \textbf{unterscheiden Sie!}\\
\end{tabular}

\vspace{10pt}

\begin{tabular}{rl}
\multicolumn{2}{l}{\textbf{PERFEKT}}\\
& haben + \textbf{unterschieden}\\
\end{tabular}
\hfill
\begin{tabular}{rl}
\multicolumn{2}{l}{\textbf{FUTURE I}}\\
& werden + \textbf{unterscheiden}\\
\end{tabular}
\hfill
\begin{tabular}{rl}
\multicolumn{2}{l}{\textbf{PLUSQUAMPERFEKT}}\\
& hatten + \textbf{unterschieden}\\
\end{tabular}

\vspace{10pt}

\begin{tabular}{lll}
\multicolumn{3}{l}{\textbf{MIT PRÄPOSITIONEN}}\\
& \textbf{unterscheiden zwischen} + Dat & (to distinguish between)\\
\end{tabular}
\hfill
\begin{tabular}{lll}
\multicolumn{3}{l}{\textbf{TRENNBARE VERBEN}}\\
& \textbf{---} & \\
\end{tabular}

\vspace{-5pt}
\end{verbentry}

% unterschreiben (to sign something)
\begin{verbentry}{
  style = {acc},
  toc = {unterschreiben (to sign something)},
  name = {unterschreiben},
  english = {to sign something},
  type = {irregular, + Akkusativ},
  auxiliary = {haben}
}
 
    
     \vspace{5pt}

    \hrule
    
    \vspace{5pt}

  \begin{tabular}{rl}
     \multicolumn{2}{l}{\textbf{PRÄSENS} }\\
    ich & \textbf{unterschreibe}\\
    du & \textbf{unterschreibst}\\
    er/sie/es & \textbf{unterschreibt}\\
    wir & \textbf{unterschreiben}\\
    ihr & \textbf{unterschreibt}\\
    sie/Sie & \textbf{unterschreiben}\\
  \end{tabular}
  \hfill
     \begin{tabular}{rl}
    \multicolumn{2}{l}{\textbf{PRÄTERITUM} }\\
    ich & \textbf{unterschrieb}\\
    du & \textbf{unterschriebst}\\
    er/sie/es & \textbf{unterschreib}\\
    wir & \textbf{unterschrieben}\\
    ihr & \textbf{unterschriebt}\\
    sie/Sie & \textbf{unterschrieben}\\
  \end{tabular}
  \hfill
  \begin{tabular}{rl}
     \multicolumn{2}{l}{\textbf{IMPERATIV} }\\
   & \textbf{unterschreib(e) (du)!}\\
   & \textbf{unterschreiben wir!}\\
   & \textbf{unterschreibt (ihr)!}\\
   & \textbf{unterschreiben Sie!}\\
   & \vspace{10pt}
  \end{tabular}

    \vspace{10pt}

 \begin{tabular}{rl}
     \multicolumn{2}{l}{\textbf{PERFEKT}}\\
   & haben + \textbf{unterschrieben}\\
  \end{tabular}
\hfill
   \begin{tabular}{rl}
   
    \multicolumn{2}{l}{\textbf{FUTURE I} }\\
   & werden + \textbf{unterschreiben}\\
 
  \end{tabular}
\hfill
   \begin{tabular}{rl}
      \multicolumn{2}{l}{\textbf{PLUSQUAMPERFEKT}}\\
   & hatten + \textbf{unterschrieben}\\
  \end{tabular}
  
  
\vspace{10pt}

\begin{tabular}{lll}
\multicolumn{3}{l}{\textbf{MIT PRÄPOSITIONEN}}\\
& \textbf{---} & \\
\end{tabular}
\hfill
\begin{tabular}{lll}
\multicolumn{3}{l}{\textbf{TRENNBARE VERBEN}}\\
& \textbf{---} & \\
\end{tabular}
\vspace{-5pt}
\end{verbentry}

% unterstützen (to support)
\begin{verbentry}{
  style = {default},
  toc = {unterstützen (to support)},
  name = {unterstützen},
  english = {to support},
  type = {regular},
  auxiliary = {haben}
}
 
    
     \vspace{5pt}

    \hrule
    
    \vspace{5pt}

  \begin{tabular}{rl}
     \multicolumn{2}{l}{\textbf{PRÄSENS} }\\
    ich & \textbf{unterstütze}\\
    du & \textbf{unterstützt}\\
    er/sie/es & \textbf{unterstützt}\\
    wir & \textbf{unterstützen}\\
    ihr & \textbf{unterstützt}\\
    sie/Sie & \textbf{unterstützen}\\
  \end{tabular}
  \hfill
     \begin{tabular}{rl}
    \multicolumn{2}{l}{\textbf{PRÄTERITUM} }\\
    ich & \textbf{unterstützte}\\
    du & \textbf{unterstütztest}\\
    er/sie/es & \textbf{unterstützte}\\
    wir & \textbf{unterstützten}\\
    ihr & \textbf{unterstütztet}\\
    sie/Sie & \textbf{unterstützten}\\
  \end{tabular}
  \hfill
  \begin{tabular}{rl}
     \multicolumn{2}{l}{\textbf{IMPERATIV} }\\
   & \textbf{unterstütz(e) (du)!}\\
   & \textbf{unterstützen wir!}\\
   & \textbf{unterstützt (ihr)!}\\
   & \textbf{unterstützen Sie!}\\
   & \vspace{10pt}
  \end{tabular}

    \vspace{10pt}

 \begin{tabular}{rl}
     \multicolumn{2}{l}{\textbf{PERFEKT}}\\
   & haben + \textbf{unterstützt}\\
  \end{tabular}
\hfill
   \begin{tabular}{rl}
   
    \multicolumn{2}{l}{\textbf{FUTURE I} }\\
   & werden + \textbf{unterstützen}\\
 
  \end{tabular}
\hfill
   \begin{tabular}{rl}
      \multicolumn{2}{l}{\textbf{PLUSQUAMPERFEKT}}\\
   & hatten + \textbf{unterstützt}\\
  \end{tabular}

   \vspace{10pt}

   \begin{tabular}{lll}
      \multicolumn{3}{l}{\textbf{MIT PRÄPOSITIONEN}}\\
& \textbf{unterstützen bei} + Dat & (to support with / help with)\\
& \textbf{unterstützen mit} + Dat & (to support with / by means of)\\
\end{tabular}
  \hfill
   \begin{tabular}{lll}
\multicolumn{3}{l}{\textbf{TRENNBARE VERBEN}}\\
& \textbf{---} & \\
\end{tabular}


  \vspace{-5pt}
\end{verbentry}

% untersuchen (to examine)
\begin{verbentry}{
  style = {default},
  toc = {untersuchen (to examine)},
  name = {untersuchen},
  english = {to examine},
  type = {irregular},
  auxiliary = {haben}
}
 
    
     \vspace{5pt}

    \hrule
    
    \vspace{5pt}

  \begin{tabular}{rl}
     \multicolumn{2}{l}{\textbf{PRÄSENS} }\\
    ich & \textbf{untersuche}\\
    du & \textbf{untersuchst}\\
    er/sie/es & \textbf{untersucht}\\
    wir & \textbf{untersuchen}\\
    ihr & \textbf{untersucht}\\
    sie/Sie & \textbf{untersuchen}\\
  \end{tabular}
  \hfill
     \begin{tabular}{rl}
    \multicolumn{2}{l}{\textbf{PRÄTERITUM} }\\
    ich & \textbf{untersuchte}\\
    du & \textbf{untersuchtest}\\
    er/sie/es & \textbf{untersuchte}\\
    wir & \textbf{untersuchten}\\
    ihr & \textbf{untersuchtet}\\
    sie/Sie & \textbf{untersuchten}\\
  \end{tabular}
  \hfill
  \begin{tabular}{rl}
     \multicolumn{2}{l}{\textbf{IMPERATIV} }\\
   & \textbf{untersuch(e) (du)!}\\
   & \textbf{untersuchen wir!}\\
   & \textbf{untersucht (ihr)!}\\
   & \textbf{untersuchen Sie!}\\
   & \vspace{10pt}
  \end{tabular}

    \vspace{10pt}

 \begin{tabular}{rl}
     \multicolumn{2}{l}{\textbf{PERFEKT}}\\
   & haben + \textbf{untersucht}\\
  \end{tabular}
\hfill
   \begin{tabular}{rl}
   
    \multicolumn{2}{l}{\textbf{FUTURE I} }\\
   & werden + \textbf{untersuchen}\\
 
  \end{tabular}
\hfill
   \begin{tabular}{rl}
      \multicolumn{2}{l}{\textbf{PLUSQUAMPERFEKT}}\\
   & hatten + \textbf{untersucht}\\
  \end{tabular}
  
  \vspace{10pt}



\begin{tabular}{lll}
\multicolumn{3}{l}{\textbf{MIT PRÄPOSITIONEN}}\\
& \textbf{---} & \\
\end{tabular}
\hfill
\begin{tabular}{lll}
\multicolumn{3}{l}{\textbf{TRENNBARE VERBEN}}\\
& \textbf{---} & \\
\end{tabular}

  \vspace{-5pt}
\end{verbentry}
